\subsection{4.75: noncentral involutions modulo the odd-order radical}
\label{cand:4.75}
\problemstatement{4.75}
Use the same group \(G\) as in Section~\ref{cand:10.62}; this is a
second consequence of one geometric construction.
Every \(2\)-subgroup has exponent two because \(G\) has exponent \(2p\).
It is therefore abelian, and the absence of a subgroup of order four
forces its order to be at most two. The maximal \(2\)-subgroups are
exactly the cyclic groups generated by involutions, and are locally
cyclic. No infinite-group Sylow theorem is required.

For a periodic group, \(O_{2'}(G)\) denotes the largest normal subgroup
all of whose elements have odd order. It exists: if \(A,B\) are normal
odd-order-element subgroups, any involution of \(AB\) has trivial image
in \(AB/A\) and would lie in \(A\), a contradiction.
Periodicity then makes all elements of \(AB\) odd-order.
The directed union of finite products gives the largest such subgroup.
Write it as \(O\). Every involution \(i\) has nontrivial image in \(G/O\).

If \(iO\) were central, all conjugates of \(i\) would have the same
image. They generate \(G\), so \(G/O=\langle iO\rangle\) would be
cyclic. It is also a quotient of the perfect group \(G\), and hence
perfect, forcing it to be trivial. This contradicts \(i\notin O\).
Thus every involution has noncentral image in \(G/O_{2'}(G)\), answering
4.75 negatively even with Sylow \(2\)-subgroups of order two.
Determination of the radical itself is unnecessary.
See \artifact{research/4.75-proof.md}.
